\documentclass{rapportCS} 
% Définit le type de document. Ici, on utilise une classe personnalisée "rapportCS"
% (probablement définie ailleurs dans ton projet).

%------------ Debut Packages ----------------

\usepackage[T1]{fontenc} 
% Permet une meilleure gestion des accents et caractères spéciaux à l'impression (encodage des fontes en T1)

\usepackage[utf8]{inputenc} 
% Définit l'encodage des fichiers source en UTF-8 (utile pour les caractères accentués)

\usepackage{lmodern} 
% Utilise la police Latin Modern (amélioration de la police Computer Modern par défaut)

\usepackage[french]{babel} 
% Adapte le document au français (traduction automatique de mots comme "Chapter" → "Chapitre", gestion des césures...)

\usepackage{xcolor} 
% Permet d'utiliser les couleurs dans le document (texte, fond, tableaux, etc.)

\usepackage{multirow} 
% Pour fusionner plusieurs lignes dans un tableau

\usepackage{makecell} 
% Pour créer des cellules de tableaux personnalisées (ex : texte sur plusieurs lignes dans une cellule)

\usepackage{float} 
% Permet un placement plus précis des figures et tableaux (avec [H] pour figer à l’endroit exact)

\usepackage{tablefootnote} 
% Permet d’ajouter des notes de bas de page dans les tableaux

\usepackage{hyperref} 
% Pour insérer des liens cliquables (URL, références internes, etc.)

\usepackage{bookmark} 
% Améliore la gestion des signets PDF (utile en complément de hyperref)

\usepackage{pdfpages} 
% Permet d'inclure des pages PDF externes dans ton document

\usepackage[linesnumbered,ruled,vlined]{algorithm2e} 
% Pour écrire des algorithmes formatés (avec numérotation des lignes, règles et indentations)

\usepackage{lipsum} 
% Pour générer du faux texte (texte de remplissage, utile en phase de test ou maquette)

\usepackage{listings} 
% Pour insérer du code source avec mise en forme (coloration syntaxique, encadré, etc.)

\usepackage{subcaption} 
% Pour insérer plusieurs figures ou tableaux côte à côte avec des sous-titres (subfigures/subtables)

\usepackage{lscape} 
% Permet de faire pivoter certaines pages en mode paysage (utile pour les grands tableaux ou schémas)

\usepackage{url} 
% Pour gérer l’affichage des URL (utile pour que les adresses web soient bien formatées)

%------------ Fin des packages classiques ----------------

%------------ Entêtes et pieds de page ----------------

\usepackage{fancyhdr} 
% Permet de personnaliser les entêtes et pieds de pages

\pagestyle{fancy} 
% Active le style fancy pour toutes les pages standards du document

\renewcommand{\chaptermark}[1]{\markboth{\chaptername \ \thechapter.\ #1}{}} 
% Personnalise l'entête de gauche (nom du chapitre) : "Chapitre 1. Titre", sans tout mettre en majuscule

\renewcommand{\sectionmark}[1]{\markright{\thesection.\ #1}} 
% Personnalise l'entête de droite (nom de la section courante) : "1.1 Titre de section", sans majuscules

\renewcommand{\thesection}{\arabic{section}} 
% Modifie le style de numérotation des sections (exclut le numéro de chapitre, utile dans les rapports courts)

%------------ Fin Packages ---------------- 
% Importe les configurations et packages (polices, marges, langue, etc.) définis dans le fichier "preambule.tex"


\title{Rapport Stage EMSI} %Thanks for Rapport CentraleSupelec - Template, By Axel Poupart-Lafarge

\begin{document} 
%----------- Informations du rapport ---------
\logoentreprise{logos/logo-entreprise.png}

\titre{Conception et Développement d'une Plateforme de Messagerie Multi-Canal Intelligente (MessageForge)}

\mention{Ingénierie Informatique et Réseaux\\ 4\textsuperscript{ème} année (4DDSI)}
\filiere{Projet de Fin d'Année (PFA)}
\master{École Marocaine des Sciences de l'Ingénieur}

\eleve{
    Abdelaadim Admi\\
}

\dates{2025 - 2026}

% Informations tuteurs écoles
\tuteuruniv{
    \textsc{Pr. Houssam Baza}
} 

\vspace{1.5cm}
% Informations membres du jury
\membresjury{
    \textsc{Pr. Houssam Baza} (Encadrant)
} 

%----------- Initialisation -------------------
        
\fairemarges %Afficher les marges
\fairepagedegarde %Créer la page de garde


%----------- Abstract -------------------
\vspace*{\stretch{1}}
\begin{center}
	\begin{abstract}
	
        Ce rapport présente la conception et le développement de \textbf{MessageForge}, une plateforme de messagerie multi-canal intelligente et hautement extensible. L'objectif principal de ce projet est de permettre aux utilisateurs d'écrire un message unique sous forme de brouillon brut et de le distribuer automatiquement et simultanément sur plusieurs canaux de communication (Email, SMS, Twitter, Slack, WhatsApp, Facebook, LinkedIn, Telegram), tout en respectant les contraintes strictes de chaque canal (limites de caractères, formats requis, enrichissements spécifiques). L'application s'appuie sur des architectures orientées services et des patrons de conception éprouvés (\textit{Strategy, Decorator, Factory, Template Method, Observer, Builder}). Développé avec le framework \textbf{Spring Boot 3.3.0}, PostgreSQL, Redis et RabbitMQ, le backend a été dockerisé et optimisé pour le déploiement local et la production.
        
        \rule{\linewidth}{0.2 mm} \\[0.4 cm]
       
        \begin{center}\textbf{Abstract}\end{center} 
        
        This report details the design and implementation of \textbf{MessageForge}, a smart and highly extensible multi-channel messaging platform. The primary goal of this project is to allow users to write a single raw draft and automatically format and dispatch it across various communication channels (Email, SMS, Twitter, Slack, WhatsApp, Facebook, LinkedIn, Telegram), fully complying with each channel's specific formatting rules and character constraints. The system leverages robust software engineering design patterns (\textit{Strategy, Decorator, Factory, Template Method, Observer, Builder}). Built using the \textbf{Spring Boot 3.3.0} ecosystem, PostgreSQL, Redis, and RabbitMQ, the backend has been fully containerized using Docker Compose for reliable local development and deployment.
    \end{abstract}
\end{center}
\vspace*{\stretch{1}}


%----------- Remerciements -------------------
\chapter*{Remerciements}
\markboth{Remerciements}{Remerciements}

Je tiens tout d'abord à exprimer ma profonde gratitude à l'administration de l'**École Marocaine des Sciences de l'Ingénieur (EMSI)** pour la qualité de la formation théorique et pratique reçue tout au long de mon parcours universitaire.

Mes remerciements s'adressent tout particulièrement à mon encadrant pédagogique, **Pr. Houssam Baza**, pour ses conseils précieux, sa disponibilité continue, ses remarques constructives et son soutien indéfectible tout au long de la réalisation de ce Projet de Fin d'Année (PFA).

Je tiens également à remercier chaleureusement tous les membres du jury qui ont accepté d'évaluer ce travail et de participer à ma soutenance. Leurs commentaires et observations contribueront sans aucun doute à enrichir et parfaire la qualité de ce projet.

Enfin, je remercie ma famille et mes collègues qui m'ont encouragé et soutenu moralement durant la phase de réalisation de ce projet.

%------------ Table des matières ----------------
\tabledematieres % Créer la table de matières

%------------ Corps du rapport ----------------
%------------ Introduction ----------------
\chapter*{Introduction Générale}
\label{chapter:introduction}
\addcontentsline{toc}{chapter}{Introduction Générale}

Dans l'ère numérique actuelle, la communication multi-canal est devenue un pilier fondamental pour la relation client, le marketing et la collaboration interne au sein des organisations. Les entreprises doivent diffuser leurs messages sur une multitude de plateformes telles que les e-mails professionnels, les SMS mobiles, les canaux collaboratifs (Slack, Microsoft Teams) ainsi que les réseaux sociaux (Twitter/X, Facebook Messenger, LinkedIn, WhatsApp, Telegram).

Cependant, chaque canal impose ses propres contraintes techniques et structurelles. Par exemple, un SMS est limité à 160 caractères et nécessite de masquer les URLs complexes pour éviter d'être bloqué par les opérateurs. Twitter/X impose une limite de 280 caractères, tandis que Slack supporte un format de balisage Markdown personnalisé. À l'opposé, l'e-mail requiert des modèles HTML riches et sécurisés. Réécrire manuellement chaque message pour s'adapter à ces spécifications est un processus chronophage, sujet aux erreurs de saisie et inefficace à grande échelle.

Pour résoudre cette problématique, nous avons conçu et développé **MessageForge**, une plateforme de messagerie intelligente et centralisée. L'application permet à un utilisateur de saisir un message brut unique et d'automatiser sa décoration (génération d'emojis, raccourcissement d'URLs, intégration de hashtags et signatures) ainsi que son formatage pour chaque canal de destination, avant de l'expédier de manière asynchrone et fiable.

Ce rapport détaille les phases clés du développement de MessageForge et est structuré comme suit :
\begin{itemize}
    \item \textbf{Chapitre 1 : Étude bibliographique et état de l'art} — Une revue des technologies clés telles que Spring Boot, GraalVM Spring Native, et les middlewares de file d'attente (RabbitMQ) et de cache (Redis).
    \item \textbf{Chapitre 2 : Analyse des besoins et spécifications} — Définition des exigences fonctionnelles et non-fonctionnelles de la plateforme.
    \item \textbf{Chapitre 3 : Architecture et conception technique} — Présentation détaillée des patrons de conception (\textit{Design Patterns}) implémentés pour garantir la modularité, ainsi que la modélisation de la base de données.
    \item \textbf{Chapitre 4 : Réalisation, tests et validation} — Détails de l'implémentation, de la dockerisation complète du système et des tests d'intégration réalisés via requêtes API.
\end{itemize}


%------------- Chapitres ----------------
\chapter{Étude bibliographique et État de l'art}
\label{chapter:etat_de_l_art}

Ce chapitre présente une étude approfondie des concepts technologiques clés et des outils logiciels utilisés pour concevoir la plateforme de messagerie MessageForge. Nous analysons l'évolution des architectures Java modernes vers les conteneurs et la compilation native, ainsi que le rôle des bases de données et des middlewares dans les systèmes distribués.

\section{Spring Boot et Spring Native (GraalVM)}

Traditionnellement, les applications Java s'exécutent sur une Machine Virtuelle Java (JVM) en utilisant la compilation à la volée (\textit{Just-In-Time} ou JIT). Bien que performante, cette approche souffre d'un temps de démarrage lent (souvent plusieurs secondes) et d'une consommation mémoire élevée lors de la phase de préchauffage (\textit{warmup}), ce qui est pénalisant dans les environnements de conteneurs légers (Docker, Kubernetes) ou serverless \cite{springnative}.

Pour surmonter ces limites, le projet **Spring Native**, en s'appuyant sur **GraalVM**, introduit la compilation en amont (\textit{Ahead-Of-Time} ou AOT) \cite{graalvm}. 
\begin{itemize}
    \item \textbf{Fonctionnement} : L'outil d'analyse statique de GraalVM examine le bytecode de l'application depuis son point d'entrée, élimine le code mort (classes, méthodes et champs inutilisés) et compile le reste directement en un fichier exécutable natif propre au système d'exploitation cible.
    \item \textbf{Avantages} : Les exécutables obtenus démarrent instantanément (en quelques dizaines de millisecondes) et consomment beaucoup moins de mémoire vive (RAM) dès le démarrage, tout en éliminant le besoin d'installer un JRE sur la machine hôte.
\end{itemize}

\section{PostgreSQL et le stockage flexible JSONB}

Pour MessageForge, le choix s'est porté sur **PostgreSQL 16**. Ce système de gestion de base de données relationnelle (SGBDR) offre à la fois la robustesse transactionnelle (ACID) pour la gestion des utilisateurs et des messages de base, ainsi que la flexibilité du format non-relationnel **JSONB** (JSON Binaire).
Le format JSONB stocke les documents sous forme décomposée en mémoire, permettant d'indexer efficacement les champs imbriqués et de réaliser des requêtes ultra-rapides. Dans notre projet, le type JSONB est utilisé pour stocker :
\begin{itemize}
    \item Les configurations personnalisables des intégrations des utilisateurs (clés API, configurations spécifiques).
    \item Les métadonnées enrichies des messages et les listes de décorateurs appliqués à la volée.
\end{itemize}

\section{Optimisation des performances avec Redis}

**Redis** (Remote Dictionary Server) est une base de données de structure de données en mémoire, utilisée comme cache et magasin de clés-valeurs. 
Dans les architectures multi-canal, les requêtes répétitives sur les configurations d'intégration ou les jetons de session des utilisateurs peuvent surcharger la base de données relationnelle. L'intégration de Redis en amont de PostgreSQL permet de mettre en cache les profils utilisateurs et les jetons de rafraîchissement (Refresh Tokens), garantissant un temps d'accès inférieur à la milliseconde et améliorant la réactivité globale des APIs.

\section{Communication asynchrone avec RabbitMQ}

Lorsqu'un message doit être envoyé à plusieurs canaux (par exemple, envoyer un e-mail via Mailgun et un SMS via Twilio), le temps d'attente lié aux appels réseaux vers les APIs externes est imprévisible et potentiellement long. Si l'application traitait ces envois de manière synchrone dans le thread de la requête HTTP, l'utilisateur ferait face à de longs temps de latence et le serveur risquerait la saturation.

Pour résoudre ce problème, un agent de messages (\textit{Message Broker}) comme **RabbitMQ** est inséré dans l'architecture \cite{rabbitmqconf}. 
\begin{itemize}
    \item \textbf{Principe} : Lorsqu'un utilisateur demande l'envoi d'un message, l'API REST crée les tâches correspondantes sous forme de messages AMQP, les publie dans un échangeur (Exchange) RabbitMQ et retourne immédiatement un accusé de réception à l'utilisateur (le message passe à l'état \texttt{SENDING}).
    \item \textbf{Traitement Asynchrone} : En arrière-plan, des consommateurs (\textit{workers}) récupèrent les tâches depuis les files d'attente (Queues) RabbitMQ de manière asynchrone, appellent les APIs externes et mettent à jour le statut final dans la base de données. En cas de panne d'une API tierce, RabbitMQ permet de stocker temporairement les tâches dans des files d'attente de lettres mortes (\textit{Dead Letter Queues}) pour retenter l'envoi ultérieurement, sans perte de données.
\end{itemize}

\chapter{Spécifications et Analyse des Besoins}
\label{chapter:specifications}

Ce chapitre détaille l'analyse des besoins de la plateforme MessageForge. Nous y définissons les exigences fonctionnelles requises pour répondre aux attentes des utilisateurs, ainsi que les exigences non-fonctionnelles garantissant la sécurité, la performance et la robustesse de l'application.

\section{Exigences Fonctionnelles}

Les exigences fonctionnelles décrivent les actions spécifiques que le système doit être capable d'exécuter. Pour MessageForge, elles se divisent en cinq grands modules :

\subsection{1. Gestion de l'Authentification et des Sessions}
Le système doit assurer une gestion sécurisée des sessions utilisateurs :
\begin{itemize}
    \item \textbf{Inscription} : Permettre à un nouvel utilisateur de s'enregistrer en saisissant son nom, son adresse e-mail et son mot de passe.
    \item \textbf{Connexion} : Authentifier l'utilisateur et générer un jeton d'accès (\textit{Access Token}) à courte durée de vie et un jeton de rafraîchissement (\textit{Refresh Token}) à longue durée de vie.
    \item \textbf{Rafraîchissement} : Permettre au client d'obtenir un nouvel Access Token à l'aide du Refresh Token sans se reconnecter.
    \item \textbf{Profil} : Récupérer les informations de profil de l'utilisateur connecté de manière sécurisée.
\end{itemize}

\subsection{2. Gestion des Brouillons de Messages (CRUD)}
L'utilisateur doit pouvoir gérer ses messages bruts sous forme de brouillons :
\begin{itemize}
    \item Créer un message brut avec un titre et un contenu textuel libre (comprenant des URLs, des hashtags ou des mentions).
    \item Consulter, modifier ou supprimer un message existant (uniquement si le message est encore à l'état de brouillon ou d'échec d'envoi).
    \item Dupliquer un message existant pour créer rapidement une nouvelle version de travail.
\end{itemize}

\subsection{3. Gestion des Intégrations de Canaux}
L'utilisateur doit pouvoir lier et configurer ses propres accès aux APIs externes :
\begin{itemize}
    \item Activer et configurer un canal spécifique (Email, SMS, Twitter, Slack, etc.) en saisissant les informations d'authentification requises (ex: jeton Slack, clé API Twilio, mot de passe SMTP).
    \item Désactiver temporairement une intégration sans supprimer la configuration sous-jacente.
    \item Déclencher un test de connectivité unitaire (\textit{Test Integration}) pour vérifier que les identifiants saisis sont valides.
\end{itemize}

\subsection{4. Moteur d'Aperçu de Formatage Dynamique}
Avant d'envoyer le message, le système doit générer des aperçus formatés en temps réel :
\begin{itemize}
    \item Pour chaque canal actif de l'utilisateur, appliquer les décorateurs sélectionnés (raccourcisseur de lien, signatures, hashtags, emojis).
    \item Adapter le texte final selon les limites de caractères du canal (ex: insérer des points de suspension à la fin si Twitter dépasse 280 caractères, ou remplacer les liens par \texttt{[URL]} pour le SMS).
    \item Retourner le nombre de caractères finaux, le statut de troncature et d'éventuels avertissements.
\end{itemize}

\subsection{5. Envoi Asynchrone des Messages}
\begin{itemize}
    \item L'utilisateur peut déclencher l'envoi d'un message vers une liste de canaux cibles.
    \item L'envoi doit s'effectuer de manière asynchrone en arrière-plan pour éviter de bloquer l'interface.
    \item Le système doit proposer un \textbf{mode test} permettant de simuler l'envoi vers les APIs externes tout en générant le flux d'événements complet.
\end{itemize}

\section{Exigences Non-Fonctionnelles}

Les exigences non-fonctionnelles décrivent les caractéristiques de qualité du système (sécurité, performance, évolutivité) :

\begin{enumerate}
    \item \textbf{Sécurité des Données} : Les mots de passe doivent être hachés en base de données à l'aide de l'algorithme BCrypt. De plus, les jetons API et clés d'intégration saisis par les utilisateurs doivent être chiffrés de manière forte dans la base PostgreSQL en utilisant le protocole AES-256-GCM.
    \item \textbf{Performance et Temps de Réponse} : L'application doit démarrer en moins de 15 secondes au sein de son conteneur Docker et répondre aux requêtes de santé (Actuator) en moins de 500ms.
    \item \textbf{Robustesse et Tolérance aux Pannes} : Une panne sur une API tierce (ex: panne temporaire de l'API LinkedIn) ne doit pas bloquer ou faire échouer les envois sur les autres canaux (SMS ou E-mail). Les tâches d'envoi doivent être isolées de manière unitaire.
    \item \textbf{Limitation de Débit (Rate Limiting)} : Pour prévenir les abus et les attaques de type déni de service (DoS), chaque utilisateur est limité à un maximum de 100 requêtes API par minute. Les requêtes excédentaires doivent recevoir un code d'erreur HTTP 429 (\textit{Too Many Requests}).
    \item \textbf{Extensibilité de l'Architecture} : Le code doit être structuré de manière à pouvoir ajouter un 9\textsuperscript{ème} canal (par exemple, Discord) ou un nouveau décorateur de texte en écrivant une classe isolée, sans modifier le cœur de l'application.
\end{enumerate}

\section{Cas d'Utilisation Principal}

Le diagramme de cas d'utilisation ci-dessous résume les interactions de l'utilisateur connecté avec MessageForge :

\begin{table}[H]
    \centering
    \caption{Cas d'utilisation : Interactions utilisateur}
    \label{tab:use_cases}
    \begin{tabular}{|l|p{10cm}|}
        \hline
        \textbf{Acteur} & Utilisateur Authentifié \\
        \hline
        \textbf{Cas d'utilisation} & 
        1. S'authentifier (Connexion / Rafraîchir Session) \newline
        2. Configurer et tester des intégrations (SMS, Twitter, Slack...) \newline
        3. Rédiger et éditer un brouillon de message \newline
        4. Demander un aperçu en direct formaté par canal \newline
        5. Lancer l'envoi d'un message vers des canaux sélectionnés \\
        \hline
    \end{tabular}
\end{table}
\chapter{Architecture et Conception Technique}
\label{chapter:conception}

Ce chapitre présente la conception détaillée de l'application. Nous y décrivons l'implémentation rigoureuse des patrons de conception (\textit{Design Patterns}), l'architecture des WebSockets pour la communication en temps réel et la modélisation de la base de données.

\section{Patrons de Conception Implémentés}

Pour garantir la flexibilité, la réutilisabilité et la conformité aux principes SOLID (notamment le principe Ouvert/Fermé), nous avons structuré le moteur de traitement des messages autour de plusieurs patrons de conception :

\subsection{1. Patron Strategy (Formatage des Canaux)}
Le formatage d'un message brut diffère grandement selon le canal cible. Pour éviter des structures conditionnelles complexes (\texttt{if-else}), nous avons implémenté le patron \textbf{Strategy} \cite{designpatterns}.
\begin{itemize}
    \item \textbf{Interface} : L'interface \texttt{MessageFormatterStrategy} déclare deux méthodes clés : \texttt{format(String rawContent, JsonNode settings)} et \texttt{validate(String content)}.
    \item \textbf{Stratégies Concrètes} : Chaque canal (Email, SMS, Twitter, Slack, LinkedIn, WhatsApp, Facebook, Telegram) implémente cette interface dans une classe dédiée (ex: \texttt{SmsFormatterStrategy}, \texttt{TwitterFormatterStrategy}).
    \item \textbf{Contexte} : La classe \texttt{FormatterContext} maintient une table de hachage de toutes les stratégies disponibles. Au moment de générer les aperçus ou d'envoyer un message, le contexte sélectionne dynamiquement la bonne stratégie en fonction du type de canal.
\end{itemize}

\subsection{2. Patron Decorator (Enrichissement de Messages)}
Avant le formatage final par canal, le message brut peut être enrichi d'éléments optionnels (emojis, signatures, hashtags, raccourcissement d'URLs). Ces modifications doivent pouvoir être combinées de façon dynamique. Le patron \textbf{Decorator} répond parfaitement à ce besoin \cite{designpatterns}.
\begin{itemize}
    \item \textbf{Structure} : La classe abstraite \texttt{MessageDecorator} implémente \texttt{MessageFormatterStrategy} et contient une référence vers une autre instance de \texttt{MessageFormatterStrategy} (le composant enveloppé).
    \item \textbf{Décorateurs Concrets} :
    \begin{itemize}
        \item \texttt{EmojiDecorator} : Ajoute automatiquement des emojis pertinents.
        \item \texttt{HashtagDecorator} : Génère ou ajoute des hashtags en fin de message.
        \item \texttt{SignatureDecorator} : Ajoute la signature de l'expéditeur.
        \item \texttt{UrlShortenerDecorator} : Détecte les liens HTTP et les remplace par des URLs courtes sous le domaine de la plateforme.
    \end{itemize}
    \item \textbf{Intérêt} : Les décorateurs peuvent être chaînés dans n'importe quel ordre à l'exécution (ex: appliquer d'abord le raccourcisseur d'URL, puis la signature, puis les hashtags).
\end{itemize}

\subsection{3. Patron Observer (Notifications Réseau)}
Pour notifier l'utilisateur en temps réel de l'état d'avancement des envois sans coupler le service d'envoi et les contrôleurs réseau, nous utilisons le patron \textbf{Observer} sous forme d'événements Spring.
Le service d'envoi (\texttt{MessageSenderService}) publie des événements de type \texttt{MessageStatusEvent} via l'interface \texttt{ApplicationEventPublisher}. Le contrôleur de WebSocket (\texttt{WebSocketController}) écoute ces événements et transmet les mises à jour aux clients abonnés.

\subsection{4. Patron Builder (Construction d'Objets)}
Les entités complexes telles que \texttt{Message} et \texttt{ChannelMessage} comportent de nombreux champs optionnels. L'utilisation du patron \textbf{Builder} (généré par l'annotation Lombok \texttt{@Builder}) nous permet d'instancier ces objets de manière claire, lisible et immuable.

\section{Architecture Temps Réel via WebSockets (STOMP)}

Pour offrir une expérience interactive, MessageForge utilise le protocole **STOMP** (\textit{Simple Text Oriented Messaging Protocol}) sur des canaux WebSocket. 

La classe \texttt{WebSocketConfig} configure un point de connexion sur le chemin \texttt{/ws-preview} et active un courtier de messages simple (\textit{Simple Message Broker}) gérant les flux suivants :
\begin{itemize}
    \item **Abonnement** : Le client s'abonne à la destination \texttt{/topic/preview.\{userId\}} pour recevoir ses aperçus de messages, et à \texttt{/topic/send.progress.\{userId\}} pour suivre la progression de l'envoi en direct.
    \item **Échange** : Le client envoie ses requêtes d'aperçu rapide à \texttt{/app/preview.request}. Le contrôleur traite la requête et renvoie immédiatement l'aperçu formaté sur le topic de l'utilisateur.
\end{itemize}

\section{Modélisation de la Base de Données}

Le modèle de données relationnel de MessageForge est conçu pour être à la fois structuré pour la cohérence transactionnelle et flexible pour les métadonnées variables.

Le diagramme de la base de données s'articule autour des tables principales suivantes :
\begin{enumerate}
    \item \textbf{users} : Enregistre les comptes utilisateurs (identifiant UUID, e-mail unique, nom d'affichage et mot de passe haché).
    \item \textbf{messages} : Stocke les brouillons de messages créés par les utilisateurs, avec leur contenu brut, leur titre et un champ \texttt{metadata} de type JSONB.
    \item \textbf{channel\_messages} : Contient le contenu mis en forme et finalisé pour chaque canal cible, lié au message principal par une clé étrangère (\textit{Foreign Key}).
    \item \textbf{channel\_integrations} : Enregistre les configurations d'accès pour chaque canal (ex: clés API Twilio, tokens Slack) associées à chaque utilisateur. Les données sensibles (identifiants d'API) y sont stockées de façon chiffrée.
    \item \textbf{refresh\_tokens} : Gère le cycle de vie des jetons de session (jetons JWT de rafraîchissement) pour le maintien des connexions.
\end{enumerate}

Le tableau ci-dessous décrit les relations physiques entre ces tables :

\begin{table}[H]
    \centering
    \caption{Description des Clés Étrangères et Relations}
    \label{tab:database_relations}
    \begin{tabular}{|l|l|l|l|}
        \hline
        \textbf{Table Source} & \textbf{Champ Source} & \textbf{Table Cible} & \textbf{Type de Relation} \\
        \hline
        \texttt{messages} & \texttt{user\_id} & \texttt{users} & Plusieurs-à-Un (N:1) \\
        \hline
        \texttt{channel\_messages} & \texttt{message\_id} & \texttt{messages} & Plusieurs-à-Un (N:1) \\
        \hline
        \texttt{channel\_integrations} & \texttt{user\_id} & \texttt{users} & Plusieurs-à-Un (N:1) \\
        \hline
        \texttt{refresh\_tokens} & \texttt{user\_id} & \texttt{users} & Un-à-Un (1:1) \\
        \hline
    \end{tabular}
\end{table}
\chapter{Réalisation, Tests et Validation}
\label{chapter:realisation}

Ce chapitre présente la mise en œuvre pratique de la plateforme MessageForge. Nous décrivons les choix techniques de l'environnement de développement, la configuration Docker et Docker Compose du projet, ainsi que les étapes de tests et de validation de nos APIs.

\section{Environnement de Développement}

La réalisation du projet s'est faite avec l'environnement technique suivant :
\begin{itemize}
    \item \textbf{Langage} : Java 21 (LTS).
    \item \textbf{Framework principal} : Spring Boot 3.3.0.
    \item \textbf{Gestionnaire de base de données} : PostgreSQL 16 (Image Docker officielle).
    \item \textbf{Gestionnaire de cache} : Redis 7 (Image Docker officielle).
    \item \textbf{Agent de messages} : RabbitMQ 3.13 (Image Docker officielle avec console d'administration).
    \item \textbf{Gestionnaire de dépendances} : Apache Maven 3.9.4.
\end{itemize}

\section{Dockerisation de l'Application}

Pour simplifier le déploiement local et assurer la reproductibilité de l'application, nous avons mis en œuvre une configuration Docker multi-stage stable et performante.

\subsection{1. Le Fichier Dockerfile}
La compilation est divisée en deux étapes distinctes (Build et Runtime) afin de minimiser la taille de l'image finale et de renforcer la sécurité en excluant le compilateur et le code source de l'environnement de production.

Le code du fichier \texttt{Dockerfile} est présenté ci-dessous :

\begin{lstlisting}[caption={Dockerfile multi-stage de l'application}, label=lst:dockerfile, frame=single, basicstyle=\footnotesize\ttfamily, keepspaces=true]
# Stage 1: Build standard executable JAR
FROM maven:3.9.4-eclipse-temurin-21 AS builder
WORKDIR /app
COPY pom.xml .
COPY src ./src
RUN mvn clean package -DskipTests

# Stage 2: Runtime JRE
FROM eclipse-temurin:21-jre-jammy
WORKDIR /app
COPY --from=builder /app/target/messageforge-1.0.0.jar /app/messageforge.jar
RUN apt-get update && apt-get install -y curl && rm -rf /var/lib/apt/lists/*
HEALTHCHECK --interval=30s --timeout=3s --start-period=15s --retries=3 \
    CMD curl -f http://localhost:8080/api/actuator/health || exit 1
EXPOSE 8080
ENTRYPOINT ["java", "-jar", "/app/messageforge.jar"]
\end{lstlisting}

\subsection{2. Orchestration via Docker Compose}
Le fichier \texttt{docker-compose.yml} configure l'ensemble des dépendances nécessaires au fonctionnement du système : PostgreSQL, Redis, RabbitMQ et l'interface d'administration pgAdmin. 

L'application Java (\texttt{messageforge-app}) y est déclarée en liant ses variables de connexion aux hôtes internes de Docker Compose (\texttt{postgres}, \texttt{redis}, \texttt{rabbitmq}). Le port hôte de l'application a été redirigé sur le port **\texttt{8085}** afin de résoudre les conflits avec d'éventuels serveurs Tomcat locaux fonctionnant sur le port par défaut \texttt{8080}.

\section{Validation et Tests d'Intégration}

Nous avons validé le bon fonctionnement de l'ensemble du système à travers un scénario de test complet exécuté via l'outil en ligne de commande \texttt{curl}.

\subsection{Étape A : Inscription de l'utilisateur}
Nous créons un compte de test en appelant l'API d'inscription :
\begin{lstlisting}[basicstyle=\footnotesize\ttfamily, frame=single]
curl -X POST http://127.0.0.1:8085/api/auth/register \
  -H "Content-Type: application/json" \
  -d "{\"email\":\"test@example.com\",\"password\":\"Test@1234\",\"displayName\":\"Test User\"}"
\end{lstlisting}

La base PostgreSQL enregistre le nouvel utilisateur et retourne la réponse suivante :
\begin{lstlisting}[basicstyle=\footnotesize\ttfamily, frame=single]
{
  "id": "f066cafe-6b55-417f-a796-5ea5713e43d1",
  "email": "test@example.com",
  "displayName": "Test User",
  "isActive": true,
  "createdAt": "2026-06-06T18:29:27"
}
\end{lstlisting}

\subsection{Étape B : Connexion de l'utilisateur}
Nous connectons l'utilisateur pour obtenir un jeton JWT d'accès sécurisé :
\begin{lstlisting}[basicstyle=\footnotesize\ttfamily, frame=single]
curl -X POST http://127.0.0.1:8085/api/auth/login \
  -H "Content-Type: application/json" \
  -d "{\"email\":\"test@example.com\",\"password\":\"Test@1234\"}"
\end{lstlisting}
Cette requête renvoie un jeton signé (\texttt{accessToken}) que nous injectons dans l'en-tête \texttt{Authorization: Bearer <TOKEN>} des requêtes suivantes.

\subsection{Étape C : Configuration des Intégrations}
Par défaut, le moteur de preview ne génère des résultats que pour les canaux activement configurés par l'utilisateur. Nous activons ici l'intégration **SMS** et **Twitter/X** :
\begin{lstlisting}[basicstyle=\footnotesize\ttfamily, frame=single]
# Activation de l'integration SMS
curl -X PUT http://127.0.0.1:8085/api/integrations/SMS \
  -H "Content-Type: application/json" \
  -H "Authorization: Bearer <TOKEN>" \
  -d "{\"enabled\":true,\"credentials\":{},\"settings\":{}}"
\end{lstlisting}

\subsection{Étape D : Création d'un Brouillon de Message}
Nous rédigeons un message brut contenant une URL et un hashtag pour tester les filtres de décoration et de formatage automatique :
\begin{lstlisting}[basicstyle=\footnotesize\ttfamily, frame=single]
curl -X POST http://127.0.0.1:8085/api/messages \
  -H "Content-Type: application/json" \
  -H "Authorization: Bearer <TOKEN>" \
  -d "{\"title\":\"First Message\",\"rawContent\":\"Hello World! Please check out https://google.com/deepmind for more details #welcome\"}"
\end{lstlisting}
Le système génère un brouillon et nous renvoie son identifiant unique, par exemple \texttt{875fb73d-3412-4231-9f65-f593809e0c23}.

\subsection{Étape E : Génération et Validation des Aperçus}
Nous appelons l'API d'aperçu dynamique pour valider la mise en conformité du texte selon les canaux :
\begin{lstlisting}[basicstyle=\footnotesize\ttfamily, frame=single]
curl -X GET http://127.0.0.1:8085/api/messages/875fb73d-3412-4231-9f65-f593809e0c23/preview \
  -H "Authorization: Bearer <TOKEN>"
\end{lstlisting}

Le serveur retourne les aperçus formatés en appliquant les règles métiers définies au chapitre 3 :
\begin{lstlisting}[basicstyle=\footnotesize\ttfamily, frame=single]
[
  {
    "channelType": "SMS",
    "formattedContent": "Hello World! Please check out [URL] for more details #welcome",
    "characterCount": 61,
    "characterLimit": 160,
    "requiresTruncation": false
  },
  {
    "channelType": "TWITTER",
    "formattedContent": "Hello World! Please check out https://google.com/deepmind for more details #welcome",
    "characterCount": 83,
    "characterLimit": 280,
    "requiresTruncation": false
  }
]
\end{lstlisting}

**Analyse du Résultat** : 
Le moteur a correctement traité les messages de façon différenciée : l'URL a été masquée par le jeton \texttt{[URL]} pour le SMS afin d'éviter le blocage antispam, alors qu'elle a été conservée au complet pour Twitter/X, confirmant la bonne exécution en chaîne des patrons Strategy et Decorator.

%------------- Conclusion Générale ----------------
\chapter*{Conclusion Générale}
\label{chapter:Conclusion}
\addcontentsline{toc}{chapter}{Conclusion Générale}
\setcounter{section}{0}

\section{Conclusion}

Ce Projet de Fin d'Année (PFA) a permis de concevoir et de réaliser la plateforme de messagerie multi-canal intelligente \textbf{MessageForge}. Nous avons réussi à mettre en œuvre un backend complet et robuste en s'appuyant sur le framework Spring Boot 21, PostgreSQL, Redis et RabbitMQ, permettant la distribution automatique, simultanée et asynchrone de messages.

L'apport majeur de ce travail réside dans la modélisation logicielle rigoureuse. L'utilisation conjointe des patrons de conception **Strategy** et **Decorator** a permis de construire un pipeline de traitement de texte modulaire, capable de raccourcir les liens, injecter des emojis ou des signatures, et valider la longueur finale du texte selon le canal cible de manière dynamique. L'intégration des WebSockets (STOMP) offre une dimension interactive en temps réel, essentielle pour les éditeurs de contenu moderne.

Ce projet a également été l'opportunité de maîtriser les processus modernes de conteneurisation et d'orchestration (Docker et Docker Compose) et de se confronter aux problématiques de gestion de ressources et de performance mémoire dans les environnements de compilation native.

\section{Perspectives}

Plusieurs évolutions peuvent être envisagées pour étendre les fonctionnalités de MessageForge :

\begin{itemize}
    \item \textbf{Intégration de l'Intelligence Artificielle} : L'ajout d'un module d'IA générative (via des LLMs comme GPT ou Gemini) permettrait de reformuler dynamiquement le contenu du message brut pour adapter automatiquement le ton (professionnel pour LinkedIn, synthétique pour Twitter, convivial pour Slack) au lieu de se limiter à des règles statiques.
    \item \textbf{Développement d'une Interface Frontend} : Concevoir un client web dynamique en utilisant des frameworks modernes tels que Next.js ou React, afin d'exploiter pleinement le flux de prévisualisation en temps réel via WebSockets.
    \item \textbf{Tableau de Bord Analytique} : Mettre en œuvre un module de statistiques pour analyser les taux de livraison, les temps de réponse des APIs externes et l'engagement des utilisateurs par canal.
\end{itemize}


\bibliographystyle{ieeetr}
\bibliography{bibliography}

\end{document} 
% Fin du document